%2345678901234567890123456789012345678901234567890123456789012345678901234567890
%        1         2         3         4         5         6         7         8

\documentclass[conference,letterpaper,10pt]{IEEEtran} 

\usepackage[T1]{fontenc}
\usepackage[utf8]{inputenc}
\usepackage[spanish,es-noshorthands]{babel}
\usepackage{lmodern}
\usepackage{microtype}
\usepackage{amsmath}
\usepackage{booktabs}
\usepackage{graphicx}
\usepackage{float}
\usepackage{xcolor}
\usepackage{pgfplots}
\pgfplotsset{compat=1.18}

\definecolor{azulUNAM}{HTML}{182C50}
\definecolor{turquesa}{HTML}{168D87}
\definecolor{grisGrafica}{HTML}{82909F}
\definecolor{grisTexto}{HTML}{536173}
\usepackage{graphicx}

\title{\LARGE \bf
Reporte de rendimiento\\
Ray Tracer
}

\author{\IEEEauthorblockN{Equipo Compumagos\\Alan Alvarez, Pedro Pablo, Fernanda Miranda}
\IEEEauthorblockA{Facultad de Ciencias, Universidad Nacional Autónoma de México\\
Modelado y Programación}}

\begin{document}

\maketitle
\thispagestyle{empty}
\pagestyle{empty}

\section{INTRODUCCIÓN}
Se desarrollo un proyecto de Ray tracer, encargado de leer, interpretar, renderizar y escribir
imagenes PPM (P3). Este reporte explicara y mostrara el rendimineto
respecto a una escena \texttt{escena\_rendimiento.json}, de 8 objetos, a 400 x 225 píxeles, tiempo de renderizado con 1,2,4 y 8 trabajadores.
Veremos el SpeedUp y la interpretacion de ley de Amdahl.
El tiempo se midió con \texttt{System.nanoTime()} alrededor
de la llamada \texttt{renderizar}. Por tanto, incluye el cálculo
y la coordinación de hilos, pero excluye la lectura del JSON
y la escritura del PPM.


\section{RESULTADOS}

\vspace{-0.3em}
\begin{table}[!ht]
\centering
\caption{Tiempos de renderizado en segundos y aceleración.}
\label{tab:resultados}
\scriptsize
\setlength{\tabcolsep}{3.4pt}
\renewcommand{\arraystretch}{1.20}
\begin{tabular}{@{}c r r r r r@{}}
\toprule
\(Trabajadores\) & \(T_1\) & \(T_2\) & \(T_3\) & \(\overline{T}_p\) & \(S(p)\) \\
\midrule
1 & 6.906 & 6.844 & 6.935 & \textbf{6.895} & 1.00 \\
2 & 3.768 & 3.572 & 3.682 & \textbf{3.674} & 1.88 \\
4 & 2.094 & 2.079 & 2.145 & \textbf{2.106} & 3.27 \\
8 & 1.524 & 1.463 & 1.691 & \textbf{1.559} & 4.42 \\
\bottomrule
\end{tabular}
\end{table}
\vspace{-1.0em}

Las métricas se obtuvieron mediante:
\begin{equation}
\overline{T}_p=\frac{1}{3}\sum_{k=1}^{3}T_{p,k}, \qquad
S(p)=\frac{\overline{T}_1}{\overline{T}_p}.
\label{eq:speedup}
\end{equation} 

\(\overline{T}_p\): Tiempo promedio.
\begin{itemize}
  \item \(p\): número de trabajadores (\(1, 2, 4\) u \(8\)).
  \item \(k\): número de ejecución (\(1, 2\) o \(3\)).
  \item \(T_{p,k}\): tiempo obtenido utilizando \(p\) trabajadores en la ejecución \(k\).\\
\end{itemize}
\(S(p)  \): Aceleracion obtenida con \(p\) trabajadores
\begin{itemize}
  \item \(\overline{T}_1\): tiempo promedio con un solo trabajador.
  \item \(\overline{T}_p\): tiempo promedio con \(p\) trabajadores.\\
\end{itemize}

Con ocho trabajadores, el tiempo promedio pasó de
\textbf{6.895 s} a \textbf{1.559 s}, una reducción aproximada del \textbf{77.4\,\%}.\\
Es decir, el renderizado es aproximadamente \textbf{4.42 veces más rápido}.


\section{ANÁLISIS DE ACELERACIÓN}
\vspace{-0.4em}
\begin{figure}[H]
\centering
\begin{tikzpicture}
\begin{axis}[
    width=\columnwidth,
    height=4.45cm,
    xmin=0.7, xmax=8.4,
    ymin=0, ymax=8.5,
    xtick={1,2,4,8},
    ytick={0,2,4,6,8},
    xlabel={Número de trabajadores},
    ylabel={Aceleración},
    xlabel style={font=\footnotesize},
    ylabel style={font=\footnotesize},
    tick label style={font=\scriptsize, color=grisTexto},
    axis line style={color=grisGrafica},
    grid=major,
    grid style={draw=gray!18},
    legend style={font=\scriptsize,draw=none,fill=none,at={(0.03,0.98)},anchor=north west},
    legend cell align=left,
    clip=false
]
\addplot[gray!65,dashed,very thick,mark=none] coordinates {(1,1) (2,2) (4,4) (8,8)};
\addlegendentry{Ideal: \(S(p)=p\)}
\addplot[turquesa,very thick,mark=*,mark size=2.5pt] coordinates {(1,1) (2,1.8767) (4,3.2740) (8,4.4218)};
\addlegendentry{Speedup observado}
\end{axis}
\end{tikzpicture}
\vspace{-0.5em}
\label{fig:speedup}
\end{figure}
\vspace{-1.0em}
La gráfica compara la aceleración ideal con la obtenida experimentalmente. 
Podemos observar que, al aumentar el número de trabajadores, el tiempo de renderizado disminuye y la aceleración aumenta.
Sin embargo, nuestros resultados se alejan progresivamente de la aceleración ideal. 
Por ejemplo, con ocho trabajadores obtuvimos una aceleración de 4.42 veces, en lugar de ocho.

\section{LEY DE AMDAHL}
La ley de Amdahl establece que, bajo un modelo de condiciones perfectas y que la cantidad de trabajo es constante
sin contar los trabajadores que usemos, la aceleracion se limita por la fraccion secuencial.
\begin{equation}
S_{\mathrm{Amdahl}}(p)=\frac{1}{(1-f)+f/p},
\label{eq:amdahl}
\end{equation}

donde $f$ representa la fracción paralelizable del programa, $1-f$ la fracción secuencial y $p$ el número de trabajadores.

Como observamos en la gráfica, con ocho trabajadores obtuvimos una aceleración de $4.42\times$, en lugar de las $8\times$ que
se alcanzarían en condiciones ideales.

Esto ocurre porque no todo el trabajo puede paralelizarse. Además, la creación de hilos, la sincronización mediante
\texttt{join()} y la planificación del procesador también pueden afectar el rendimiento. No medimos estos costos por
separado, por lo que no podemos determinar cuánto contribuye cada uno a la diferencia observada.


\section{CONCLUCIÓN}

El paralelismo por filas redujo de forma notoria el tiempo de renderizado. No obstante, la ganancia fue
decreciente: la eficiencia $S(p)/p$ pasó de \textbf{81.8\,\%}  con cuatro
trabajadores a \textbf{55.3\,\%} con ocho, en concordancia con los
límites del paralelismo.

\addtolength{\textheight}{-12cm} 
\end{document}
